\documentclass[10pt,a4paper]{amsart}
\usepackage{amsmath,amssymb,amsthm,graphicx}
% Retained minimal preamble for the Stacks Project chapter more-algebra.tex (More on Algebra).
% The chapter file pinned-original-source.tex carries only \input{preamble} on its first line, so no
% preamble text lives inside the chapter itself.  The Stacks Project uses one single shared preamble for
% every chapter; the snapshot stages it identically as sources/stacks-chow/preamble.tex and
% sources/stacks-etale-cohomology/preamble.tex.  The block below is copied mechanically, line for line,
% from sources/stacks-chow/preamble.tex lines 14-18 (the Xy-pic setup) and 154-217 (the native theorem
% declarations together with the mathematical macros).  Nothing was added, removed or redefined; the
% wrapper class and the amsmath/amssymb/amsthm/graphicx base packages are supplied by the assembler.
% Chapter: more-algebra.tex -- More on Algebra.
% For commutative diagrams we use Xy-pic
\usepackage[all]{xy}

% We use 2cell for 2-commutative diagrams.
\xyoption{2cell}

% Theorem environments.
%
\theoremstyle{plain}
\newtheorem{theorem}[subsection]{Theorem}
\newtheorem{proposition}[subsection]{Proposition}
\newtheorem{lemma}[subsection]{Lemma}

\theoremstyle{definition}
\newtheorem{definition}[subsection]{Definition}
\newtheorem{example}[subsection]{Example}
\newtheorem{exercise}[subsection]{Exercise}
\newtheorem{situation}[subsection]{Situation}

\theoremstyle{remark}
\newtheorem{remark}[subsection]{Remark}
\newtheorem{remarks}[subsection]{Remarks}

\numberwithin{equation}{subsection}

% Macros
%
\def\lim{\mathop{\mathrm{lim}}\nolimits}
\def\colim{\mathop{\mathrm{colim}}\nolimits}
\def\Spec{\mathop{\mathrm{Spec}}}
\def\Hom{\mathop{\mathrm{Hom}}\nolimits}
\def\Ext{\mathop{\mathrm{Ext}}\nolimits}
\def\SheafHom{\mathop{\mathcal{H}\!\mathit{om}}\nolimits}
\def\SheafExt{\mathop{\mathcal{E}\!\mathit{xt}}\nolimits}
\def\Sch{\mathit{Sch}}
\def\Mor{\mathop{\mathrm{Mor}}\nolimits}
\def\Ob{\mathop{\mathrm{Ob}}\nolimits}
\def\Sh{\mathop{\mathit{Sh}}\nolimits}
\def\NL{\mathop{N\!L}\nolimits}
\def\CH{\mathop{\mathrm{CH}}\nolimits}
\def\proetale{{pro\text{-}\acute{e}tale}}
\def\etale{{\acute{e}tale}}
\def\QCoh{\mathit{QCoh}}
\def\Ker{\mathop{\mathrm{Ker}}}
\def\Im{\mathop{\mathrm{Im}}}
\def\Coker{\mathop{\mathrm{Coker}}}
\def\Coim{\mathop{\mathrm{Coim}}}

% Boxtimes
%
\DeclareMathSymbol{\boxtimes}{\mathbin}{AMSa}{"02}

%
% Macros for moduli stacks/spaces
%
\def\QCohstack{\mathcal{QC}\!\mathit{oh}}
\def\Cohstack{\mathcal{C}\!\mathit{oh}}
\def\Spacesstack{\mathcal{S}\!\mathit{paces}}
\def\Quotfunctor{\mathrm{Quot}}
\def\Hilbfunctor{\mathrm{Hilb}}
\def\Curvesstack{\mathcal{C}\!\mathit{urves}}
\def\Polarizedstack{\mathcal{P}\!\mathit{olarized}}
\def\Complexesstack{\mathcal{C}\!\mathit{omplexes}}
% \Pic is the operator that assigns to X its picard group, usage \Pic(X)
% \Picardstack_{X/B} denotes the Picard stack of X over B
% \Picardfunctor_{X/B} denotes the Picard functor of X over B
\def\Pic{\mathop{\mathrm{Pic}}\nolimits}
\def\Picardstack{\mathcal{P}\!\mathit{ic}}
\def\Picardfunctor{\mathrm{Pic}}
\def\Deformationcategory{\mathcal{D}\!\mathit{ef}}

% --- appended 2026-10-12 (Chief): the project own definitions of the reference/slogan/history
% display environments, copied mechanically from the shared Stacks preamble lines 8-12
% (\usepackage{verbatim} plus the three \newenvironment lines, which the chapter preamble keeps
% verbatim from the project preamble). Nothing added, removed or redefined: these environments
% carry no mathematical content and the project itself typesets them as comments. Without these
% three lines a counted statement whose own text uses them fails with Environment slogan undefined.
\usepackage{verbatim}
\newenvironment{reference}{\comment}{\endcomment}
\newenvironment{slogan}{\comment}{\endcomment}
\newenvironment{history}{\comment}{\endcomment}
\usepackage{cleveref}
\crefname{originallemma}{original lemma}{original lemmas}
\Crefname{originallemma}{Original lemma}{Original lemmas}
\makeatletter
\newlabel{algebra-lemma-Artin-Rees}{{00IN}{source}{}{sourcepointer-1}{}}
\newlabel{algebra-lemma-Artin-Rees@cref}{{[originallemma][0][]00IN}{[1][0][]source}{}{}{}}
\makeatother
\crefname{originallemma}{original lemma}{original lemmas}
\Crefname{originallemma}{Original lemma}{Original lemmas}
\makeatletter
\newlabel{derived-lemma-essentially-constant-cohomology}{{0G3B}{source}{}{sourcepointer-2}{}}
\newlabel{derived-lemma-essentially-constant-cohomology@cref}{{[originallemma][0][]0G3B}{[1][0][]source}{}{}{}}
\makeatother
\crefname{originallemma}{original lemma}{original lemmas}
\Crefname{originallemma}{Original lemma}{Original lemmas}
\makeatletter
\newlabel{derived-lemma-pro-isomorphism}{{0G3C}{source}{}{sourcepointer-3}{}}
\newlabel{derived-lemma-pro-isomorphism@cref}{{[originallemma][0][]0G3C}{[1][0][]source}{}{}{}}
\makeatother
\crefname{originallemma}{original lemma}{original lemmas}
\Crefname{originallemma}{Original lemma}{Original lemmas}
\makeatletter
\newlabel{derived-remark-truncation-distinguished-triangle}{{08J5}{source}{}{sourcepointer-4}{}}
\newlabel{derived-remark-truncation-distinguished-triangle@cref}{{[originallemma][0][]08J5}{[1][0][]source}{}{}{}}
\makeatother
\crefname{originallemma}{original lemma}{original lemmas}
\Crefname{originallemma}{Original lemma}{Original lemmas}
\makeatletter
\newlabel{lemma-homotopy-koszul}{{6660}{source}{}{sourcepointer-5}{}}
\newlabel{lemma-homotopy-koszul@cref}{{[originallemma][0][]6660}{[1][0][]source}{}{}{}}
\makeatother
\crefname{originaldefinition}{original definition}{original definitions}
\Crefname{originaldefinition}{Original definition}{Original definitions}
\makeatletter
\newlabel{situation-koszul}{{26305}{source}{}{sourcepointer-6}{}}
\newlabel{situation-koszul@cref}{{[originaldefinition][0][]26305}{[1][0][]source}{}{}{}}
\makeatother
\title{For a Noetherian ring the Koszul complexes of a sequence are pro-isomorphic to the quotients by the powers of the sequence}
\author{The Stacks Project Authors}
\date{stacks-project:more-algebra.tex:local-snapshot-20261012; GNU FDL 1.2 or later}
\begin{document}
\maketitle
\setcounter{subsection}{1}
\noindent\textit{Editorial note.} All delivered mathematical text is the author's own, copied from the pinned Stacks Project chapter file more-algebra.tex; no statement or proof step was rewritten, reordered or removed.\par\medskip
\noindent\textit{Editorial note.} The chapter file pinned-original-source.tex carries only \textbackslash{}input\{preamble\} on its first line, so no preamble text lives inside the chapter. The retained file sources/stacks-more-algebra/retained-preamble.tex is a mechanical, line-for-line copy of the project's single shared preamble as staged in the snapshot; nothing was added, removed or redefined and no package was added.\par\medskip
\noindent\textit{Editorial note.} The source URL is the source-level Stacks Project URL, as required when more than one card is produced from one source; the per-claim tag is carried by claim\_key.\par\medskip
\noindent\textit{Editorial note.} The counted statement span includes the author's own footnote explaining the consequence of the pro-isomorphism for the factorizations; the footnote is delivered verbatim. The six references inside the counted proof span are recorded as non-load-bearing navigation pointers.\par\medskip
\noindent\textit{Allowed standard prerequisites (editorial).}\begin{itemize}
\item The criterion for an inverse system of distinguished triangles to induce a pro-isomorphism. invoked at proof lines 27188-27190. Reduces the pro-isomorphism statement about K\_n\textasciicircum{}bullet to pro-zeroness of the truncation system.
\item The essentially constant cohomology criterion for a bounded complex. invoked at proof lines 27192-27195. Reduces pro-zeroness of the truncation system to pro-zeroness of the individual cohomology systems.
\item The Koszul complex of a sequence, its degree bounds and the annihilation of its cohomology by the ideal generated by the powers. used at proof lines 27191-27192 and 27206-27207. Supplies the degree range and the annihilation that the transition-map computation needs.
\item Artin-Rees applied to an ideal and a submodule of a finite module. invoked at proof lines 27209-27211. Is the input that forces the transition map on cohomology to vanish.
\end{itemize}\medskip
\ifdefined\hypertarget\hypertarget{sourcepointer-1}{}\fi
\noindent\textit{Source reference.} Original lemma 00IN: focused-seeds/AGP/sources/stacks-snapshot-20261011/tags.txt, source lines 685--685. Cross-chapter navigation pointer into the Stacks Project; the standard fact it supplies is applied inside the counted proof span and the argument resting on it stays there.\par\medskip
\ifdefined\hypertarget\hypertarget{sourcepointer-2}{}\fi
\noindent\textit{Source reference.} Original lemma 0G3B: focused-seeds/AGP/sources/stacks-snapshot-20261011/tags.txt, source lines 20086--20086. Cross-chapter navigation pointer into the Stacks Project; the standard fact it supplies is applied inside the counted proof span and the argument resting on it stays there.\par\medskip
\ifdefined\hypertarget\hypertarget{sourcepointer-3}{}\fi
\noindent\textit{Source reference.} Original lemma 0G3C: focused-seeds/AGP/sources/stacks-snapshot-20261011/tags.txt, source lines 20087--20087. Cross-chapter navigation pointer into the Stacks Project; the standard fact it supplies is applied inside the counted proof span and the argument resting on it stays there.\par\medskip
\ifdefined\hypertarget\hypertarget{sourcepointer-4}{}\fi
\noindent\textit{Source reference.} Original lemma 08J5: focused-seeds/AGP/sources/stacks-snapshot-20261011/tags.txt, source lines 10838--10838. Cross-chapter navigation pointer into the Stacks Project; the standard fact it supplies is applied inside the counted proof span and the argument resting on it stays there.\par\medskip
\ifdefined\hypertarget\hypertarget{sourcepointer-5}{}\fi
\noindent\textit{Source reference.} Original lemma 6660: sources/stacks-more-algebra/pinned-original-source.tex, source lines 6660--6660. Same-chapter navigation pointer; the statement it names is standard for the section and the argument that consumes it stays inside the counted proof span.\par\medskip
\ifdefined\hypertarget\hypertarget{sourcepointer-6}{}\fi
\noindent\textit{Source reference.} Original definition 26305: sources/stacks-more-algebra/pinned-original-source.tex, source lines 26305--26305. Same-chapter navigation pointer; the statement it names is standard for the section and the argument that consumes it stays inside the counted proof span.\par\medskip
\begin{lemma}
\label{lemma-sequence-Koszul-complexes}
In Situation \ref{situation-koszul}. If $A$ is Noetherian, then the pro-objects
$\{K_n^\bullet\}$ and $\{A/(f_1^n, \ldots, f_r^n)\}$ of $D(A)$ are
isomorphic\footnote{In particular, for every $n$ there exists an $m \geq n$
such that $K_m^\bullet \to K_n^\bullet$ factors through the map
$K_m^\bullet \to A/(f_1^m, \ldots, f_r^m)$.}.
\end{lemma}

\begin{proof}
We have an inverse system of distinguished triangles
$$
\tau_{\leq -1}K_n^\bullet \to K_n^\bullet \to A/(f_1^m, \ldots, f_r^m) \to
(\tau_{\leq -1}K_n^\bullet)[1]
$$
See Derived Categories, Remark
\ref{derived-remark-truncation-distinguished-triangle}.
By Derived Categories, Lemma \ref{derived-lemma-pro-isomorphism}
it suffices to show that the inverse system
$\tau_{\leq -1}K_n^\bullet$ is pro-zero.
Recall that $K_n^\bullet$ has nonzero terms only in degrees
$i$ with $-r \leq i \leq 0$. Thus by
Derived Categories, Lemma \ref{derived-lemma-essentially-constant-cohomology}
it suffices to show that $H^p(K_n^\bullet)$ is pro-zero
for $p \leq -1$.
In other words, for every $n \in \mathbf{N}$ we have to show
there exists an $m \geq n$ such that $H^p(K_m^\bullet) \to H^p(K_n^\bullet)$
is zero. Since $A$ is Noetherian, we see that
$$
H^p(K_n^\bullet) =
\frac{\Ker(K_n^p \to K_n^{p + 1})}{\Im(K_n^{p - 1} \to K_n^p)}
$$
is a finite $A$-module. Moreover, the map $K_m^p \to K_n^p$ is given
by a diagonal matrix whose entries are in the ideal
$(f_1^{m - n}, \ldots, f_r^{m - n})$ as $p < 0$.
Note that $H^p(K_n^\bullet)$ is annihilated by
$J = (f_1^n, \ldots, f_r^n)$, see Lemma \ref{lemma-homotopy-koszul}. Now
$(f_1^{m - n}, \ldots, f_r^{m - n}) \subset J^t$ for $m - n \geq tn$.
Thus by Algebra, Lemma \ref{algebra-lemma-Artin-Rees} (Artin-Rees) applied
to the ideal $J$ and the module $M = K_n^p$ with submodule
$N = \Ker(K_n^p \to K_n^{p + 1})$ for $m$ large enough
the image of $K_m^p \to K_n^p$ intersected with
$\Ker(K_n^p \to K_n^{p + 1})$ is contained in
$J \Ker(K_n^p \to K_n^{p + 1})$. For such $m$ we get the zero map.
\end{proof}

\end{document}
